\documentclass[UTF8,fontset=fandol,zihao=5]{ctexart}

\usepackage[a4paper,margin=2cm]{geometry}
\usepackage{amsmath,amssymb,mathtools}
\usepackage{booktabs,tabularx,array}
\usepackage{xcolor}
\usepackage[most]{tcolorbox}
\usepackage{tikz}

\usetikzlibrary{
  positioning,
  calc,
  decorations.pathreplacing,
  arrows.meta
}

\usepackage{enumitem}
\usepackage{indentfirst}
\usepackage{hyperref}

\DeclareTextFontCommand{\textbf}{\heiti\color{main}}

% relationbox 中知识关系图的统一样式
\tikzset{
  relnode/.style={
    font=\small,
    align=center
  },
  relline/.style={
    thick
  },
  relarrow/.style={
    -{Stealth},
    thick
  }
}

\definecolor{main}{HTML}{2F6B4F}
\definecolor{mainlight}{HTML}{FFFFFF}
\definecolor{method}{HTML}{4D6B5D}
\definecolor{methodlight}{HTML}{F1F5F2}

\hypersetup{colorlinks=true,linkcolor=main,urlcolor=main}
\setlength{\parindent}{2em}
\linespread{1.5}
\setlength{\abovedisplayskip}{5pt}
\setlength{\belowdisplayskip}{5pt}
\setlength{\abovedisplayshortskip}{3pt}
\setlength{\belowdisplayshortskip}{3pt}
\setlist[itemize]{leftmargin=2em,itemsep=0.1em,topsep=0.2em}
\setlist[enumerate]{leftmargin=2.3em,itemsep=0.1em,topsep=0.2em}

\ctexset{
  section={format=\centering\heiti\zihao{3}\color{main},
           name={第,章},number=\chinese{section},aftername=\quad,
           beforeskip=0.9em,afterskip=0.5em},
  subsection={format=\heiti\zihao{5}\color{main},
              name={,、},number=\chinese{subsection},aftername=,
              beforeskip=0.6em,afterskip=0.22em},
  subsubsection={format=\heiti\zihao{5}\color{main},
                 number={（\chinese{subsubsection}）},aftername=,
                 beforeskip=0.55em,afterskip=0.15em}
}

\newtcolorbox{relationbox}{
  enhanced,breakable,colback=mainlight,colframe=main,
  boxrule=0.65pt,arc=1.4mm,left=2mm,right=2mm,top=1.2mm,bottom=1.2mm
}
\newtcolorbox{methodbox}[1][]{
  enhanced,breakable,colback=methodlight,colframe=method,
  boxrule=0.65pt,arc=1.4mm,left=2mm,right=2mm,top=1.2mm,bottom=1.2mm,
  title=#1,fonttitle=\kaishu,fontupper=\kaishu,
  before upper={\setlength{\parindent}{2em}}
}

\newcommand{\jj}{\mathrm{j}}

\begin{document}

\begin{center}
  {\heiti\zihao{1}\color{main} 电路分析基础总结}
\end{center}

\vspace{0.45em}

电路分析基础研究\textbf{集总参数电路}，即电路尺寸远小于电磁波长的电路。与集总参数电路相对的是分布参数电路。电路中每个二端元件就是一个\textbf{支路}，支路的连接点叫作\textbf{节点}，电路中任一闭合路径叫作\textbf{回路}，内部不含支路的回路叫作\textbf{网孔}。此外，分析电路时还应注意\textbf{参考方向}。

\section{元件约束}

\subsection{独立源}

\begin{enumerate}
  \item 分类：电压源和电流源；理想源与实际源。

  \item VCR

  \begin{center}
    \begin{tabularx}{0.92\textwidth}{>{\bfseries}p{2.5cm}XX}
      \toprule
      & 数量形式 & 相量形式 \\
      \midrule
      电压源 & \(v=V_s\) & \(\dot V=\dot V_s\) \\
      电流源 & \(i=I_s\) & \(\dot I=\dot I_s\) \\
      \bottomrule
    \end{tabularx}
  \end{center}

  \item 功率：电阻电路中，电源的功率一般为负。

  \item 理想源等效变换：

  理想电压源串联：\(V_{\mathrm{eq}}=\displaystyle\sum V_s\)，一般不并联。\textbf{与元件并联的电压源等效为该电压源}。

  理想电流源并联：\(I_{\mathrm{eq}}=\displaystyle\sum I_s\)，一般不串联。\textbf{与元件串联的电流源等效为该电流源}。

  \item 实际源等效变换：

  \begin{center}
    \begin{tabularx}{0.92\textwidth}{XX}
      \toprule
      实际电压源：\(v=v_s-R_si\) & 实际电流源：\(i=i_s-G_sv\) \\
      \bottomrule
    \end{tabularx}
  \end{center}

  上两式满足关系
  \[
    \begin{cases}
      v_s=R_si_s,\\
      R_sG_s=1.
    \end{cases}
  \]

  \item 有源元件与无源元件

  若元件吸收的能量对任意 \(t\)、\(v\)、\(i\) 恒满足
  \[
    w=\int_{-\infty}^{t}v(\xi)i(\xi)\,\mathrm{d}\xi\geq 0,
  \]
  则该元件无源，否则有源。
\end{enumerate}

\subsection{受控源}

\begin{enumerate}
  \item 压控压源（VCVS）：\(v_2=\mu v_1\)，其中 \(\mu\) 为电压比系数。
  \item 压控流源（VCCS）：\(i_2=gv_1\)，其中 \(g\) 为转移电导。
  \item 流控流源（CCCS）：\(i_2=\alpha i_1\)，其中 \(\alpha\) 为电流比系数。
  \item 流控压源（CCVS）：\(v_2=ri_1\)，其中 \(r\) 为转移电阻。
\end{enumerate}

\subsection{线性非时变电阻}

\begin{enumerate}
  \item VCR（关联参考方向）

  \begin{center}
    \begin{tabularx}{0.92\textwidth}{XX}
      \toprule
      数量形式：\(v=Ri\quad i=Gv\) & 相量形式：\(\dot V=R\dot I\quad \dot I=G\dot V\) \\
      \bottomrule
    \end{tabularx}
  \end{center}

  \item 功率：\(p=\dfrac{v^2}{R}=i^2R\)。电阻电路中，电阻的功率恒为正。

  \item 等效电阻

  \begin{center}
    \begin{tabularx}{0.92\textwidth}{XX}
      \toprule
      串联：\(R_{\mathrm{eq}}=\displaystyle\sum R\) & 并联：\(G_{\mathrm{eq}}=\displaystyle\sum G\) \\
      \bottomrule
    \end{tabularx}
  \end{center}

  \item 输入电阻：\(R_i=\dfrac{u}{i}\)
\end{enumerate}

\subsection{电容}

\begin{enumerate}
  \item VCR（关联参考方向）

  \begin{center}
    \begin{tabularx}{0.92\textwidth}{XX}
      \toprule
      数量形式：\(i=C\dfrac{\mathrm{d}v}{\mathrm{d}t}\) & 相量形式：\(\dot I=\jj\omega C\dot V\) \\
      \bottomrule
    \end{tabularx}
  \end{center}

  \item 导纳 \(Y_C=\jj\omega C\)，容纳 \(B_C=\omega C\)
  \item 储能：\(w=\dfrac{1}{2}Cv^2\)
  \item \textbf{并联}等效电容：\(C_{\mathrm{eq}}=\displaystyle\sum C\)
\end{enumerate}

\subsection{电感}

\begin{enumerate}
  \item VCR（关联参考方向）

  \begin{center}
    \begin{tabularx}{0.92\textwidth}{XX}
      \toprule
      数量形式：\(v=L\dfrac{\mathrm{d}i}{\mathrm{d}t}\) & 相量形式：\(\dot V=\jj\omega L\dot I\) \\
      \bottomrule
    \end{tabularx}
  \end{center}

  \item 阻抗 \(Z_L=\jj\omega L\)，感抗 \(X_L=\omega L\)
  \item 储能：\(w=\dfrac{1}{2}Li^2\)
  \item \textbf{串联}等效电感：\(L_{\mathrm{eq}}=\displaystyle\sum L\)
\end{enumerate}

\subsection{耦合电感}

\begin{enumerate}
  \item 同名端为产生磁链相互加强时，电流流入一端。

  \item VCR（同名端方向一致）

  \begin{center}
    \footnotesize
    \begin{tabularx}{\linewidth}{@{}XX@{}}
      \toprule
      数量形式：
      \(\begin{cases}
        v_1=L\dfrac{\mathrm{d}i_1}{\mathrm{d}t}+M\dfrac{\mathrm{d}i_2}{\mathrm{d}t},\\[5pt]
        v_2=L\dfrac{\mathrm{d}i_2}{\mathrm{d}t}+M\dfrac{\mathrm{d}i_1}{\mathrm{d}t}
      \end{cases}\)
      & 相量形式：
      \(\begin{cases}
        \dot V_1=\jj\omega L\dot I_1+\jj\omega M\dot I_2,\\
        \dot V_2=\jj\omega L\dot I_1+\jj\omega M\dot I_2
      \end{cases}\) \\
      \bottomrule
    \end{tabularx}
  \end{center}

  \item 耦合系数 \(k=\dfrac{M}{\sqrt{L_1L_2}}\in[0,1]\)
  \item 功率 \(p=v_1i_1+v_2i_2\)
  \item 储能（同名端方向一致）\(w=\dfrac{1}{2}L_1i_1^2+\dfrac{1}{2}L_2i_2^2+Mi_1i_2\)
\end{enumerate}

\subsection{理想变压器}

\begin{enumerate}
  \item \(\widetilde S_1+\widetilde S_2=0\)，故有 \(P_1+P_2=0\)，\(Q_1+Q_2=0\)。
  \item 设 \(n=\dfrac{N_1}{N_2}\)，则同名 \(U_1=nU_2\)，异名 \(U_1=-nU_2\)。
  \item 等效阻抗 \(Z_2=n^2Z_1\)，与同名端无关。
  \item 理想变压器的线性性质对 Thevenin 等效电路也适用。注意参考方向与同名端方向
\end{enumerate}

\section{拓扑约束}

\subsection{Kirchhoff 定律}

\begin{enumerate}
  \item Kirchhoff 电流定律（KCL）：

  \begin{center}
    \begin{tabularx}{0.92\textwidth}{XX}
      \toprule
      数量形式：\(\displaystyle\sum i=0\) & 相量形式：\(\displaystyle\sum\dot I=0\) \\
      \bottomrule
    \end{tabularx}
  \end{center}

  \item Kirchhoff 电压定律（KVL）

  \begin{center}
    \begin{tabularx}{0.92\textwidth}{XX}
      \toprule
      数量形式：\(\displaystyle\sum v=0\) & 相量形式：\(\displaystyle\sum\dot V=0\) \\
      \bottomrule
    \end{tabularx}
  \end{center}
\end{enumerate}

\subsection{电阻电路分析法}

\begin{methodbox}[【方法】节点电压法]
\begin{enumerate}
    \item 文字表述：确定零电位点，设出节点电压，对其他节点运用 KCL，求解方程组。
    
    \item 矩阵表示：\(\boldsymbol{I}=\boldsymbol{YV}\)。其中：

    \(\boldsymbol{I}\) 是节点驱动电流向量，输入为正输出为负。（只考虑独立源和受控源）

    \(\boldsymbol{Y}\) 是导纳矩阵，对角线上为自导纳，即与节点相连所有导纳之和，值为正；其他为互导纳，即两节点间导纳，值为负。

    \(\boldsymbol{V}\) 是节点电压向量，与 \(\boldsymbol{I}\) 一一对应。
\end{enumerate}
\end{methodbox}

\begin{methodbox}[【方法】网孔电流法]
\begin{enumerate}[start=2]
  \item 网孔电流法（略）
\end{enumerate}
\end{methodbox}

\subsection{线性定理}

线性电路中激励与响应具有线性，即 \(\boldsymbol{r}=\boldsymbol{K}\boldsymbol{e}\)，分量形式为 \(r(e)=\displaystyle\sum K_i r(e_i)\)，由此有以下定理。

\begin{enumerate}
  \item 齐性定理：\textbf{单一激励}的线性电路中，激励与响应成正比。
  \item 叠加定理：\textbf{线性时不变电路}中，总激励等于分激励之和。
\end{enumerate}

\subsection{替代定理（唯一性定理）}

若所有支路中电压与电流均有\textbf{唯一解}，则可用理想源代替该支路。

\subsection{Thevenin 定理和 Norton 定理}

\begin{enumerate}
  \item Thevenin 定理：任何含源线性单口网络 N 都可以用一个电压源和一个电阻的串联支路替代。
  \item Norton 定理：任何含源线性单口网络 N 都可以用一个电流源和一个电阻的并联支路替代。
\end{enumerate}

注：两定理在正弦稳态电路中同样成立，但此时电阻为复值电阻。

\begin{methodbox}[【题型】Thevenin 与 Norton 等效求法]
  等效电阻法、外加电源法、短路电流法、VCR 确定法
\end{methodbox}

\subsection{最大功率传输定理}

\begin{enumerate}
  \item 线性时不变电路：当 \(R_L=R_s\) 时，电源有最大功率 \(P_{\max}=\dfrac{V_s^2}{4R_s}\)，此时效率为 \(\eta=50\%\)。
  \item 正弦稳态电路
\end{enumerate}

\subsection{一阶动态电路三要素法}

\[
  y(t)=y(\infty)+[y(0^+)-y(\infty)]\mathrm{e}^{-\frac{t}{\tau}}
\]

其中 \(y(0^+)\) 是初值，\(y(\infty)\) 是终值，\(\tau\) 是时间常数。\(\tau=RC=\dfrac{L}{R}\)，其中 \(R\) 为 Thevenin 等效电阻。

\subsection{正弦稳态电路}

\begin{enumerate}
  \item 相位规律：串联元件电流同相，并联元件电压同相。

  \item 有效值：\(V=\dfrac{\sqrt{2}}{2}V_m\)，\(I=\dfrac{\sqrt{2}}{2}I_m\)

  \item 阻抗与导纳
  \begin{itemize}
    \item 等效阻抗：\(Z=R+\jj X\)。
    
    若 \(X>0\)，则阻抗角 \(\psi_z>0\)，电压超前电流，等效为电阻与电感的串联支路；

    若 \(X<0\)，阻抗角 \(\psi_z<0\)，电流超前电压，等效为电阻与电容的串联支路。

    \item 等效导纳：\(Y=G+\jj B\)。
    
    若 \(B>0\)，则导纳角 \(\psi_y>0\)，电流超前电压，等效为电阻与电容的并联支路；

    若 \(B<0\)，则导纳角 \(\psi_y<0\)，电压超前电流，等效为电阻与电感的并联支路。
  \end{itemize}\

  \item 功率

  \begin{enumerate}[label=(\arabic*)]
    \item 有功功率（平均功率，W）：\(P=UI\cos\varphi\)
    \item 无功功率（Var）：\(Q=UI\sin\varphi\)
    \item 视在功率（VA）：\(S=\lVert\widetilde S\rVert=\sqrt{P^2+Q^2}\)
    \item 复功率（VA）：\(\widetilde S=U\dot I^*=P+\jj Q\)
  \end{enumerate}

  其中，\(\varphi\) 为压流相位差。称 \(\lambda=\cos\varphi\) 功率因数。

  功率仍然满足 \(P=\displaystyle\sum P_i\)。
\end{enumerate}

\end{document}
